Managing and exploiting a large space of useful optimizations is \system's core feature. In this section we describe key logical and physical optimizations that can be used by \system{} --- many of which have already been implemented in our prototype. %as indicated by the circled numbers (e.g., \circlednum{1}) in \autoref{fig:intro}.

\subsection{Logical Optimizations}
\label{sec:logical_opt}
\system{} has a set of logical optimizations which can be applied to the logical plan implied by a user's program. These optimizations create new plans that are logically equivalent, but may yield significant runtime and cost savings depending on factors like operation execution cost and the selectivity of filter operations.  We have implemented two logical optimizations already --- filter reordering and convert reordering --- but we plan on adding more in the future (such as splitting and coalescing LLM requests).

{\bf Filter Reordering} simply permutes the ordering of selection filters in a logical plan wherever possible. For example, if an input logical plan has the filter sequence A, B, and C, this logical optimization would yield 5 additional plans, one for each unique permutation of filters. If each of the filters has a unique selectivity, then exactly one of these permutations will minimize the number of records processed by the system. \system\ does not initially know these selectivities (though future versions may use historical data). However, \system{}'s cost optimizer is capable of estimating these selectivities after execution begins and sample execution data is collected. Thus, by considering all filter reorderings in its logical optimization stage, \system{} can find more efficient executions of user plans. 

%\mike{Let's separate the specification from the implementation concerns.  If the user writes a program that applies convert C1, and then applies filter F2 to the result, then by default the filter F2 is allowed to use any information that came about as a result of the convert operation C1.  But if the user provides depends\_on information, then we might know that F2 does not need C1's results and it can be logically reordered.  Do you think this part is clear?}

%\mike{Another concern is that we might obtain strange behavior if we have a physical optimization that coalesces different input objects for a single LLM prompt, in which case the output of processing tuple X depends on whether it was coalesced with tuple Y.  But I think this is clearly a physical concern.}
    
\noindent {\bf Convert Reordering} enables \system{} to move convert operations around the logical plan, similar to filter reordering. This can be especially beneficial when a plan has a mixture of convert operations and filter operations, because moving an expensive convert operation after a filter which does not depend on its output can save unnecessary computation. For example, in the Real Estate Search workload, the fastest plan will apply the highly-selective text-based filters before it runs the expensive image processing step for testing whether a listing is ``modern and attractive." Importantly, the reordering takes into consideration any dependencies between conversion and filter operations to guarantee the equivalence of the reordered plans. At the moment, \system{} relies on the programmer to explicitly state dependencies between convert and filter operators. In the future, we hope add the capability for \system{} to deduce this information automatically, in a process akin to acquisitional query processing \cite{madden2003acqp}.

Two plans that are logically reordered are semantically equivalent. However, if implemented naively, a logical reordering might create a new set of physical LLM prompts that yield substantively different results.  Ensuring that the physical plan faithfully reflects the logical plan semantics is important and is described in \autoref{sec:physicalopt} below.


%\st{The system gains a lot of possible optimization flexibility through these reorderings, but the user code may be ambiguous about which reorderings are safe to perform.  For example, imagine that the user describes two sequential operators, Convert-1 and Filter-2.  If Filter-2 requires output from Convert-1 in order to make a high-quality decision, it cannot be reordered to run beforehand.  But if Filter-2 does not require Convert-1's output, we can reorder it and possibly gain performance by "pushing predicates" earlier.  By default, the system assumes that every operator implementation uses {\em all} of its program-specified input data.  The user can loosen this assumption with the {\tt depends\_on} parameter and thereby open up more potential reordering optimizations.  In future work we want the system itself to figure out which fields can be safely ignored, and thus make this parameter unnecessary.}

%Note that Finally, several operations on lines 27-34 use the {\tt depends\_on} parameter (which was described in \autoref{sec:logical_opt}.  This allows the user to tell the system which input schema fields should be used (and thus, which can be ignored) when computing a result. This information allows the system more flexibility when performing logical reordering of convert operations (in this case, allowing the system to skip image processing entirely).  

\subsection{Physical Optimizations}
\label{sec:physicalopt}
A number of recent AI optimizations, such reducing text generation latency via speculative inference~\cite{liu2023online, spector2023, leviathan2023fast, chen2023accelerating}, can naturally be deployed in many contexts, including chat, general LLM API services, and \system. But some optimizations are better or are only possible because our system's declarative programming framework permits  multiple low-level implementations that are semantically equivalent to the user's goal; such optimizations are unavailable when programming AI applications via a lower-level interface (such as direct prompting). We have implemented some of these optimizations in \system\ today, while others comprise future work.


\noindent {\bf Model Selection} --- which is implemented in our prototype --- is a simple but effective example of such an optimization. A high-level AI program comprises many tasks; for example, the code in \autoref{fig:enron-email-program} entails extraction of the {\tt Email} fields from unstructured text as well as filtering the emails by content. Thanks to its declarative language, \system\ can decompose the high-level program into smaller operations and choose the most appropriate model for each. It might be fine to use a cheap, small, fast model for easy operations, and only use the expensive model for harder ones. The idea is simple, but implementing it is not: because (1) a single program can comprise many operations, (2) the exact operation decomposition can change depending on other optimizations, (3) the difficulty of an operation can change over time, and (4) model quality can fluctuate as LLM services make updates. This optimization can easily become burdensome without \system's help.

\noindent {\bf Code Synthesis} --- which is implemented in our prototype --- involves generating synthesized code to perform specific operators dynamically. In certain real-world scenarios, tasks may not require deep semantic understanding and can be efficiently handled through synthesized code. By replacing calls to LLMs with calls to synthesized functions, we can generally reduce runtime and lower costs. The SEED system \cite{chen2024seed} used this approach to generate cost-optimized data curation pipelines with ensembles of synthesized code and LLM calls. In a similar vein, the {\sc Evaporate} system~\cite{arora2023language} used synthesized code, along with some weak supervision methods~\cite{Ratner_2017}, to replace some LLM operations. \system{} analyzes a set of sample inputs for a given conversion operation and then employs an LLM to generate a function that performs the necessary conversion. This approach can streamline processes where semantic depth is unnecessary, optimizing both performance and resource utilization.

\noindent {\bf Multi-data Prompt Marshaling} --- which is implemented in our prototype --- is the problem of deciding how user operators should be decomposed into prompts. Should data objects be processed in a row- or column-centric manner? For example, when running the code at the top of \autoref{fig:examples}, computing all of the {\tt ScientificPaper} fields in a single LLM call might allow us to process the input tokens just once, while obtaining multiple output values. But if there is a filter on the {\tt title} field, and if the {\tt citation} output field is very large, then a "column"-centric approach might be better.  Furthermore, different LLMs may be able to process more or fewer examples per invocation, depending on their context window and overall effectiveness.
Since the best decision will depend on the current state of the AI program and its input data, a human engineer attempting to implement this optimization manually would need to constantly reevaluate and possibly reimplement what goes into a particular prompt. 

\noindent {\bf Input Token Reduction} --- which is implemented in our prototype --- aims to delete portions of certain operators' input data while still obtaining high-quality results. For example, it should be possible to populate the {\tt title} and {\tt authors} fields of {\tt ScientificPaper} in \autoref{fig:examples} while ignoring almost all of the input text.  In document-processing use cases, this reduction in token size --- and thus reduction in financial cost and increase in execution speed --- can be very dramatic. By observing multiple naive examples of using {\tt convert()}  to transform a PDF to a {\tt ScientificPaper}, the declarative optimizer can run automatic "experiments" to determine which regions of the input are necessary. 
In some ways, this optimization is akin to a "micro-RAG" task, choosing salient excerpts from a source object. However, unlike a full RAG system, this optimization exists only for the life of the program execution. This is possible because \system{} can learn operator properties at the schema level. In the naive chat-processing use case, which lacks schemas and a visible workload of tasks from a single program, it is not clear how this method could be implemented. 
% \cliu{ZenDB \cite{lin2024towards} extracts hierarchical structures from templated documents and executes queries using a summary-based tree search to pinpoint relevant components.}

\noindent {\bf Output Token Reduction} aims to reduce the size of LLM generation outputs without changing the application-specific quality of these outputs. SplitWise \cite{patel2023splitwise} showed that LLM inference time is generally proportional to the size of the input and output tokens, with the output tokens having a greater impact when various LLM optimizations are activated. When \system{} is asked to find large excerpts from input objects, such as populating the {\tt citations} field of {\tt ScientificPaper}, the output token size can be very large. In these cases, it may be possible to reformulate the naive operation so it gives the same answer but with dramatically smaller output. For example, the system can insert artificial tokens into the input text and then ask the LLM to use these tokens to report which text block contains the target result.  In some early test cases, this approach can reduce output sizes from thousands of tokens down to just two.

\noindent {\bf Model Cascades} is a method that many researchers have explored in the context of image processing~\cite{10.14778/3137628.3137664, anderson2019physical, Cai2015LearningCC, Sun2013DeepCN, cascadeserve}. The core idea is to construct and exploit a series of models that accomplish the same goal, ranging from  fast/low-quality models on the low-end to expensive/high-quality models on the high-end. Operator processing starts on the inexpensive model. If the model can process the input with high confidence, the system uses its output; if not, the system gives the input to a more expensive model in the sequence. This approach is not implemented in \system\ yet, but it has been very successful in other related projects.

\noindent {\bf Knowledge Distillation} methods approximate the outputs of an expensive large-parameter model by using a smaller, cheaper (and often more limited) replica model~\cite{gu2024minillm, hinton2015distilling, Gou_2021}. Methods such as the {\em teacher / student} approach have produced small models that perform similarly to much larger ones, such as Alpaca and Vicuna~\cite{alpaca, vicuna2023}. In certain cases, it may be feasible for our system to perform knowledge distillation to obtain a small and fast model on a per-operator basis.

\noindent {\bf Workload-Aware Execution Management}, enabled by \system's ability to design entire workloads of model requests, should yield opportunities in model serving and resource management. If a distributed model service permits clients to give batch-oriented optimization hints --- such as whether a set of queries will use the same model, or to implement prefill prepacking~\cite{zhao2024prepacking} --- \system\ can exploit it.  Strategically co-scheduling LLM processing with similar prompts can optimize key-value (KV) cache reuse, substantially improving the KV cache hit rate during LLM inference \cite{liu2024optimizing,ye2024cascade}. Because the system can reformulate prompts and gather statistics about output token lengths, it should be able to design batches of model requests based on similar output token lengths, thereby minimizing waiting times for requests that finish early and boosting model request throughput.  
%\mjf{Very interesting! kind of multiple query optimization.}

%It enhances control over the input token size and accurately estimates the output token size through token reduction, context batching, and question batching. This allows for adjustments in token size as needed. LLM execution can be managed at the granularity of iterations (rather than requests) \cite{yu2022orca}, and tasks can also be batched based on similar output token lengths. This 


\subsection{Choosing An Optimization}
\label{sec:choose-opt}
In order to be valuable, the system must not just hypothesize a valuable set of optimized plans, it must actually choose one that yields concrete benefits.  \system\ follows steps (marked with circled numbers) described visually at the top of \autoref{fig:intro}, and algorithmically in Algorithm~\ref{algo:optimization}.


\begin{algorithm}
\caption{Optimized Plan Selection Algorithm}
\label{algo:optimization}
\begin{algorithmic}[1]
\REQUIRE $userCode$, $userPolicy$ \COMMENT{Step \circlednum{1}}
\STATE $logicalPlans$ = generateLogicalCandidates($userCode$) \COMMENT{Step \circlednum{2}}
\STATE $sentinelPhysicalPlans$ = getPhysicalPlans($logicalPlans$, $sentinel=True$) \COMMENT{Step \circlednum{3}}
\STATE
\STATE $performanceStatisics$ = \{\}
\FOR{0...NUM\_SAMPLES}
    \STATE $input$ = getSampledInput() 
    \STATE $stats$ = runAndComputeStatistics($sentinelPhysicalPlans$, $input$)  \COMMENT{Step \circlednum{4}}
    \STATE performanceStatistics.update($stats$) \COMMENT{Step \circlednum{5}}
\ENDFOR
\STATE
\STATE $physicalCandidates$ = getPhysicalPlans($logicalPlans$, $stats=performanceStatistics$)
\STATE $reducedCandidates$ = naiveElimination($physicalCandidates$)
\STATE $frontierCandidates$ = scoreAndEliminatePlans($reducedCandidates$, $performanceStatistics$)
\STATE
\RETURN chooseBestPlan($frontierCandidates$, $userPolicy$) \COMMENT{Step \circlednum{6}}
\end{algorithmic}
\end{algorithm}

The algorithm starts with a user program and an optimization goal (such as "minimize runtime subject to F1 > 0.5").  On line 1, it generates all possible logical plans consistent with the user program, as described in \autoref{sec:logical_opt}.  On line 2, it generates a small number of "sentinel" physical plans that are consistent with these logical plans. A sentinel plan is a simple hard-coded plan that should shed information on different selectivities and the accuracy of individual non-optimized operators. We currently generate three sentinel plans: one that uses GPT-3.5 for each convert and filter, one that uses MIXTRAL 8x7B, and one that uses GPT-4. (These are depicted in the sample-based statistics collection module at the top-right of \autoref{fig:intro}.)

On lines 4-8, the algorithm runs all of the sentinel plans on a small NUM\_SAMPLES number of inputs. The goal is to collect basic statistics about the behavior of some easy-to-understand plans.

On line 10, the system hypothesizes physical plans that embody optimizations for the logical plans.  This set can be quite large.  Our prototype generates 234, 10,140, and 1,950 physical plans for the Legal Discovery, Real Estate, and Medical Schema Matching programs, respectively.  On line 11, we use a small set of rules to eliminate plans that likely do not add to the set of useful runtime options.  On line 12, the algorithm uses statistics from lines 4-8 to score the expected runtime, financial cost, and quality of all the remaining physical plans; it then throws away any plans that are not on the Pareto frontier.

Finally, on line 14, the algorithm scores each frontier candidate in light of the user's optimization goal and returns the highest-scoring option for execution.  Estimating is done using well-known relational database planner methods, while estimating cost simply requires combining published service cost data with an accurate token consumption estimate.  Measuring quality is much harder, as described in \autoref{sec:quality}. The current system simply uses the champion model approach.

The algorithm is admittedly naive. For example, it does not have a rigorous termination criterion for the sampling phase on lines 5-8, and the sentinel plans are chosen arbitrarily. But as we will see in \autoref{sec:evaluation}, the algorithm effectively identifies plans that are close to the true optimal plan within the set generated by \system{}.

% There are a number of methods, notably including the {\em teacher / student} approach. These methods have produced a number of notable models, such as Alpaca, which was generated using the self-instruct method and the Llama model~\cite{alpaca, wang2023selfinstruct, touvron2023llama}; and Vicuna, which was also based on Llama~\cite{vicuna2023}. In certain cases, it may be feasible for our system to perform per-task knowledge distillation to obtain a small and fast model.

% \srm{I think we should not just list things that we have implemented / evaluate below -- let's be generous and include all of the ongoing work, for example Zui's stuff or cascades work that Ferdi is working on. }
% {\bf TK: INSERT SOMETHING ON MODEL CASCADES? MAYBE SOMETHING ON CACHE MANAGEMENT FROM RANA/CHUNWEI/MATT?}
% \cliu{added my thoughts on the KV cache, scheduling and streaming support. Please review.}


%\st{Multi-modality is one of the key features of our system. Beyond handling diverse data modalities such as paper PDFs, images, texts, and structured datasets, we are also planning to integrate streaming data from IoT sensors and system monitoring logs. \system{}'s concise DSL simplifies the creation of dashboards and monitoring pipelines, facilitating seamless streaming data integration and management.} \textcolor{brown}{Matt: I think this either belongs in Section 5 or maybe Section 3, but this is not discussing an optimization.}

%\tim{I would mention all the additional approaches from the last paragraph as seperate items and then cite those optimizaton. Afterwards I would emphasize that the goal of PZ is to bring all the optimization under one umbrella and make them easily accessible rather that each individual technique needs to be implemented seperately}

%To deliver more relevant input context, RAG offers a promising approach by preselecting pertinent information before inputting it into the LLM. While RAG utilizes semantic similarity, it may fall short for certain tasks. For example, a basic RAG model might mistakenly retrieve the reference section when asked about the author of a paper. The similarity between the question and the text does not always ensure similarity with the answer, limiting RAG's effectiveness in practical applications. Given the prevalence of repetitive query tasks within \system{}, particularly for tasks centered on literature where the input context often shares a similar structure, we have introduced a sample-based token reduction strategy. This process begins with a set of sample papers and a user-defined question. We initially feed the entire context to LLM models to generate an answer set. Subsequently, we use fuzzy matching to locate the source of each answer within the full input context, marking it with start and end offsets. This information allows us to create a heatmap of user queries, which serves as a reference for \system{} when handling similar tasks. With a predefined token budget, \system{} can determine which specific parts of the context are necessary to answer a given question, allowing for a smart trade-off between accuracy and cost. In scenarios where the context is incomplete or insufficient, \system{} can progressively incorporate more context or revert to using the entire context if necessary. By pinpointing relevant sections, we can significantly reduce the number of tokens processed by the LLM, thereby enhancing the speed of information retrieval and processing.


%{\bf From earlier:}
%\lei{such as the AWS paper Tim co-authored which co-optimizes database retrieval and LLM inference} Recent LLM research has focused on creating new physical optimizations for AI systems that utilize LLMs \lei{Such as Ferdi's work on optimizing model serving with a cascade ensemble}. \system{} is able to leverage these optimizations --- e.g., model selection, KV-cache re-use, and learned prompts --- while also contributing new optimizations which are critical for many AI system workloads, such as bonded queries, \textcolor{blue}{code synthesis for learned filters/projections}, and token compaction.
